\section{Geometry}
The mesh in the code is defined inspired by the classical 
viewpoint of Boundary Representation (BREP)~\cite{10.1145/1499949.1500071}.
In two-dimension case,
the element is bounded by a set of edges,
and the edges are bounded by a set of vertex.
The first and most important aspect of creating a FEM or FVM code
is indexing the geometric objects.
Assuming the the following representation of global geometric object.
\begin{equation}\label{eq:globalgeometricOBJ}
		  \begin{aligned}
					 \text{Vertex} \quad &V_i, \quad i\in ID_V,\\
					 \text{Edge} \quad &E_i, \quad i\in ID_E,\\
					 \text{Cell}  \quad &C_i, \quad i\in ID_C,
		  \end{aligned}
\end{equation}
where $ID_V$, $ID_E$ and $ID_C$ are countable 
set containing all the non repeating global index for vertices, edges and faces 
respectively. We use superscript to label the subset of global index, e.g.,
$ID^{E_i}_{V} \subset ID_{V}$ represents the global index of boundary 
vertex associated with the edge $E_i$, which is a subset of global index set 
of vertex.

At the same time, assuming the following representation 
of local geometric object.
\begin{equation}\label{eq:localgeometricOBJ}
		  \begin{aligned}
					 \text{Vertex} \quad &v_i, \quad i\in id_v,\\
					 \text{Edge} \quad &e_i, \quad i\in id_e,\\
					 \text{Cell} \quad &c_i, \quad i\in id_c
		  \end{aligned}
\end{equation}
where $id_v$, $id_e$ and $id_c$ are countable set 
containing local index for vertices, edges and faces
corresponding to a given geometrical object.
We use superscript to label the target geometrical object, e.g., 
$id_v^{E_i} = \{0,1\}$ represents the local index of boundary vertex 
associated with the edge $E_i$.

We connect global index and local index with a one-to-one local to global 
mapping as
\begin{equation}
		  f: id \rightarrow ID.
\end{equation}
In implementation, we usually don't have to actually define a local to
global mapping explicitly.
Consider an array(vector) containing
global index, the array index is a natural representation for 
local index set.
With the help of array representation, we
can reverse to a global to local mapping within each target geometrical
objects.

\subsection{General Unstructured Mesh}
We start with a brief analysis of what information should a target cell
should see from its perspective.
Each polygon cell $E_i$ should have the following links to other geometric
objects as
\begin{itemize}
		  \item Adjacent polygons connected by edges;
		  \item Additional adjacent polygons connected with vertices, 
					 but not sharing an edge;
		  \item Indicident vertices;
		  \item Edges that bounds this element.
\end{itemize}
Within the local set, we can relate edges and vertices directly.
For convinence, 
each edge $j$ should has links to:
\begin{itemize}
		  \item Two polygons that sharing this edge;
		  \item Two vertices that bounds this edge.
\end{itemize}

In FEM and FVM computation, we work with a orientable mesh, where 
"clockwise" and "countercolckwise" can be defined consistently.
Hence, we can use half edge data structure to regulate direction
in a unstructured polygon mesh.
\begin{figure}[ht]
\centering
\input chapters/tikz/halfedge.tex
		  \caption{An illustration of the half-edge data structure}
		  \label{fig:halfedge}
\end{figure}
In figure~\ref{fig:halfedge}, we present an illustration of the half edge 
data structure applied to a quadrilateral mesh. For each target cell $E_i$,
it sees the vertex $V_i$ and the local link $l_i$ to the next vertex $V_i^+$.
Hence the direction of the link $l_i:V_i \rightarrow V_i^+$ is determined from the perspective of the 
target cell $C_i$.

\subsection{Reshape Cartesian Data}
The natural of a 2D cartesian data is a pair of index $(i,j)$. However,
we usually reshape the 2D data into a 1D array for a more natural 
representation in computer. We give a more general reshape function
that project N-dimensional data to a one-dimensional data. 
\begin{algorithm}[htb]
\caption{N dimensional reshape\label{alg:reshape}}
		  Consider a N-dimensional array $\{n_0\times n_i\times \ldots n_{N-1}\}$.
		  \vspace{0.3cm}

Let $I=0$ be the result, and $a=0$ be the multiplier.
		  \vspace{0.3cm}

		  Consider an input N-dimensional index $\{i_0, i_1, \ldots, i_{N-1}\}$.
		  \vspace{0.3cm}

\begin{algorithmic} [1]
		  \For{\rm rank $r= 0,1,\ldots, N-1$} 
		  \vspace{0.3cm}
		  \State $I = I + i_r * a$.
		  \vspace{0.3cm}
		  \State $a = a\times n_r$.
		  \vspace{0.3cm}
		  \EndFor
\end{algorithmic}
\end{algorithm}
We can also reverse the process reshaping 1D data back into N-dimensional 
data with modulo and integer division operation.

\subsection{Logically Rectangular Mesh}
Mesh generation is not the focus of this project, hence
we adopt a much simpler way of implementing quadrilateral mesh,
the logically rectangular way.
We generate the quadrialteral mesh by randomly distorting each 
interior vertex coordinates.
There is no affine mapping between reference square mesh and generated
mesh. The indexing of geometrical objects remain untouched through
the vertex distortion.

In our case, all the indices are represented with cartesian coordinates.
Therefore, we can setup local to global mapping of geometric objects 
easily with liear functions.
Consider a mesh with $M\times N$ cells, the totlal number of vertex
is $(M+1)\times (N+1)$ and the total number of edges is 
$M\times (N+1) + (M+1)\times N $. 
We count horizontal edges first, hence the vertical edges are counted
starting from $M\times (N+1)$.
Let local vertex and edge number being counted in the same way 
as indicated in the mesh~\ref{},
we can define the local to global mapping with local index of vertex and
global index of cell.
Consider mapping the local index $i$ of vertex $v_i$ in the cell $C_I$ as
\begin{equation}\label{eq:gtolocal}
		  ID_{V} \ni  = 
\end{equation}
and the edge local to global mapping as

\section{Assembly and dofs}
In FEM and FVM computation, we need to associate geometrical objects with
degrees of freedom.
